\chapter{Algebra and identification}\label{ch:cI1}
\chapterbrief
{$\Omega\subset\R^2$ open and connected; the ring $\ring=\C[\tau,\xi]$ acting on distributions by $\tau\mapsto\partial_t$, $\xi\mapsto\partial_x$. Classical facts used without proof: unique factorisation in $\ring$ (Gauss), the Nullstellensatz, the Chinese remainder theorem, and the fact that distributions agreeing near every point agree.}
{every identification argument reduces a question about solutions to a question about ideals and exponential-polynomials. This chapter proves those reductions in full, then the identification theorems built on them.}
{Dim.~A1--A9. They prove Chapters~3--5 of the paper: the gcd splitting, the EP test, the identification criterion, the intrinsic operator, the spectral form, the finite-data thresholds, and the decidability results.}

\begin{notebox}{Notation for the whole chapter}
$\Var(I)\subset\C^2$ is the common zero set of an ideal $I$. For a point $z=(\tau_z,\xi_z)\in\C^2$,
\[
\mathfrak m_z=(\tau-\tau_z,\ \xi-\xi_z)
\]
is its maximal ideal, and $e^{\langle z,\cdot\rangle}$ denotes $e^{\tau_zt+\xi_zx}$. $\Ann(w)=\{Q\in\ring:\ Q(\partial)w=0\}$. $\EP(\Omega)$ is the set of restrictions to $\Omega$ of finite sums $\sum_kq_k\,e^{\langle z_k,\cdot\rangle}$ with polynomials $q_k$.
\end{notebox}

%=====================================================================
\section{The algebraic toolkit}

\begin{dimcard}{A1}{C3-L2}{Splitting off the gcd}{\proved}{lem:gcd}
Let $0\ne I\subseteq\ring$ be an ideal and $G=\Gcd(I)$. Then:
\begin{enumerate}[label=(\roman*)]
\item $I=G\cdot J$ for an ideal $J$ with $\Gcd(J)=1$;
\item $\Var(J)$ is finite;
\item $(\sqrt J)^N\subseteq J$ for some $N$, and $J=\ring$ if $\Var(J)=\emptyset$.
\end{enumerate}
Two facts are proved on the way:
\begin{itemize}
\item \textbf{(V)} two nonzero coprime polynomials have finitely many common zeros;
\item \textbf{(F)} an irreducible polynomial has infinitely many zeros; hence if $\Var(\Ann w)$ is finite, some $h\in\Ann(w)$ is not divisible by a given irreducible $P_i$.
\end{itemize}
\end{dimcard}

\begin{proofidea}
Over the field $\C(\xi)$ two coprime polynomials have a B\'ezout identity; clearing denominators gives a polynomial in $\xi$ alone that vanishes at every common zero, and likewise in $\tau$. After dividing an ideal by its gcd, a generic combination of its elements is coprime to any fixed element, and (V) applies.
\end{proofidea}

\begin{proof}
\step{1}{Fact (V): coprime polynomials meet in finitely many points}
Let $f,g\ne0$ be coprime and $K=\C(\xi)$. A common factor of positive $\tau$-degree in $K[\tau]$ would give, by the Gauss lemma, a common factor in $\C[\xi][\tau]$; so $f,g$ are coprime in the principal ideal domain $K[\tau]$. Hence there are $a,b\in K[\tau]$ with
\[
af+bg=1 .
\]
Multiplying by a common denominator,
\[
a'f+b'g=r(\xi),\qquad a',b'\in\ring,\quad r\in\C[\xi]\setminus\{0\}.
\]
At a common zero $(\tau_0,\xi_0)$ the left side vanishes, so $r(\xi_0)=0$: finitely many $\xi_0$. Exchanging $\tau$ and $\xi$ gives $s(\tau)\ne0$ with $s(\tau_0)=0$. Therefore
\[
\Var(f,g)\subseteq\{r=0\}\times\{s=0\},
\]
a finite set.

\step{2}{Fact (F): an irreducible polynomial has infinitely many zeros}
If $\deg_\tau P_i\ge1$, then for all $\xi_0$ outside the finitely many zeros of the leading coefficient, $P_i(\cdot,\xi_0)$ has a root $\tau_0$. This gives infinitely many points. If $\deg_\tau P_i=0$, then $P_i$ is an irreducible polynomial in $\xi$, so $P_i=\alpha(\xi-\xi_0)$ and $\Var(P_i)=\C\times\{\xi_0\}$.

For the consequence: if every $h\in\Ann(w)$ were divisible by $P_i$, then $\Var(P_i)\subseteq\Var(\Ann w)$, an infinite set inside a finite one.

\step{3}{(i) Dividing by the gcd}
Put $J=\{f/G:\ f\in I\}$. It is closed under sums and under multiplication by $\ring$, so it is an ideal, and $I=G\cdot J$. If a prime $\pi$ divided every element of $J$, then $G\pi$ would divide every element of $I$, against the definition of $G$. So $\Gcd(J)=1$.

\step{4}{(ii) A generic combination is coprime to a fixed element}
If $J=\ring$ there is nothing to prove. Otherwise pick $0\ne f\in J$ with distinct prime factors $\pi_1,\dots,\pi_s$. By Step~3, for each $a$ there is $g_a\in J$ with $\pi_a\nmid g_a$. For $c\in\C^s$ put
\[
g_c=\sum_{b=1}^sc_b\,g_b\in J .
\]
For each $a$ the set $H_a=\{c:\ \pi_a\mid g_c\}$ is a linear subspace, and it is proper because the $a$-th unit vector is not in it. A finite union of proper subspaces does not cover $\C^s$, so some $c$ lies in no $H_a$. Then $f$ and $g_c$ are coprime, and by Step~1
\[
\Var(J)\subseteq\Var(f,g_c)\quad\text{is finite.}
\]

\step{5}{(iii) Nullstellensatz}
The polynomials vanishing on $\Var(J)$ form $\sqrt J$, which is finitely generated, so a power of it lies in $J$. If $\Var(J)=\emptyset$, then $\sqrt J=\ring$ and $1\in J$.

\whyworks{dividing by the gcd leaves an ideal with no common curve, and two curves without common component meet in finitely many points.}
\end{proof}
\uses{Gauss lemma, Nullstellensatz.}

%---------------------------------------------------------------------
\begin{dimcard}{A2}{C3-L3}{Exponential-polynomials are the finite-dimensional functions}{\proved}{lem:ep}
For $w\in\mathcal D'(\Omega)$, $\Omega$ connected, the following are equivalent:
\[
w\in\EP(\Omega)\iff\dimc_\C\ring w<\infty\iff\Var(\Ann w)\ \text{is finite.}
\]
Moreover, with distinct $z_k$:
\begin{enumerate}[label=(\alph*)]
\item $Q(\partial)\big(q\,e^{\langle z,\cdot\rangle}\big)=\big(Q(z+\partial)q\big)\,e^{\langle z,\cdot\rangle}$ for every polynomial $q$;
\item if $\sum_kq_k\,e^{\langle z_k,\cdot\rangle}=0$ on a nonempty open set, every $q_k=0$;
\item if $w=\sum_kq_k\,e^{\langle z_k,\cdot\rangle}$ with $q_k\ne0$, then $\Var(\Ann w)=\{z_k\}$.
\end{enumerate}
\end{dimcard}

\begin{proofidea}
Exponentials are independent, so an exponential-polynomial lives on finitely many points. Conversely, a finite-dimensional derivative-stable space gives two one-variable ODEs, $p(\partial_t)w=0$ and $q(\partial_x)w=0$. On a rectangle they force $w$ to be an exponential-polynomial, and connectedness glues the rectangles together.
\end{proofidea}

\begin{proof}
\step{1}{(a) Derivatives of an exponential times a polynomial}
For one derivative,
\[
\partial_t\big(q\,e^{\langle z,\cdot\rangle}\big)=\big(\partial_tq+\tau_zq\big)\,e^{\langle z,\cdot\rangle},
\]
and the same with $\partial_x$ and $\xi_z$. Iterating and using linearity gives (a).

\step{2}{(b) Exponentials are linearly independent}
The sum is real-analytic on $\R^2$ and vanishes on an open set, so it vanishes everywhere. Suppose some $q_k\ne0$, and let $K'$ be the set of such $k$. Choose $v\in\R^2$ with $\langle z_k-z_l,v\rangle\ne0$ for $k\ne l$, and $y_0$ with $\prod_{K'}q_k(y_0)\ne0$. Restricting to the line $y_0+sv$ gives
\[
f(s)=\sum_{k\in K'}r_k(s)\,e^{\mu_ks}\equiv0,\qquad \mu_k=\langle z_k,v\rangle\ \text{distinct},\quad r_k(0)\ne0 .
\]
We remove one term at a time. Multiply by $e^{-\mu_1s}$ and differentiate $\deg r_1+1$ times. The first term disappears. Each other term becomes
\[
\Big(\big(D+\mu_k-\mu_1\big)^{\deg r_1+1}r_k\Big)\,e^{(\mu_k-\mu_1)s},
\]
and since $\mu_k\ne\mu_1$ the operator $(D+c)^N$ with $c\ne0$ preserves the degree, so the new coefficient is nonzero. Induction ends with $r(s)e^{\mu s}\equiv0$, $r\ne0$: impossible.

\step{3}{(c) The zero set of the annihilator is the set of exponents}
By (a), the derivatives of $q\,e^{\langle z,\cdot\rangle}$ stay in $\C[t,x]_{\le\deg q}\cdot e^{\langle z,\cdot\rangle}$, which is finite-dimensional. If $|\beta|>\deg q$, the product $(\tau-\tau_z)^{\beta_1}(\xi-\xi_z)^{\beta_2}$ acts on $q$ as $\partial^\beta$ and kills it. So
\[
\prod_k\mathfrak m_{z_k}^{\deg q_k+1}\subseteq\Ann(w),\qquad \Var(\Ann w)\subseteq\{z_k\}.
\]
Conversely, if $Q\in\Ann(w)$ then $Q(z_k+\partial)q_k=0$ by (a) and (b). The top-degree part of $Q(z_k+\partial)q_k$ is $Q(z_k)$ times the top-degree part of $q_k$, so $Q(z_k)=0$.

\step{4}{Finite zero set implies finite dimension}
If $\Var(\Ann w)=\emptyset$ then $\Ann w=\ring$ and $w=0$. Otherwise, by A1(iii),
\[
\Ann(w)\supseteq\mathfrak I(\Var)^N,
\]
where $\mathfrak I(\Var)$ is the ideal of the finite set $\Var=\Var(\Ann w)$. By the Chinese remainder theorem
\[
\ring/\mathfrak I(\Var)^N\cong\bigoplus_{z\in\Var}\ring/\mathfrak m_z^N,
\]
which is finite-dimensional. Hence $\ring w\cong\ring/\Ann(w)$ is finite-dimensional too.

\step{5}{Finite dimension gives two ODEs}
Let $W=\ring w$, stable under $\partial_t$ and $\partial_x$. Let $p$ and $q$ be the minimal polynomials of $\partial_t|_W$ and $\partial_x|_W$. Then
\[
p(\partial_t)\,w=0,\qquad q(\partial_x)\,w=0 .
\]

\step{6}{On a rectangle the ODE in $t$ separates the variables}
Let $Q_0=I'\times J'\subset\Omega$ be an open rectangle. We claim that every $v\in\mathcal D'(Q_0)$ with $p(\partial_t)v=0$ has the form
\[
v=\sum_{j,l}c_{jl}\otimes t^le^{\lambda_jt},\qquad c_{jl}\in\mathcal D'(J'),
\]
with $\lambda_j$ the roots of $p$. Induction on $\deg p$. Write $p=(\tau-\lambda)\tilde p$. By induction $(\partial_t-\lambda)v=\sum\tilde c_{jl}\otimes t^le^{\lambda_jt}$. Each summand has a particular solution $\tilde c\otimes g(t)$:
\[
g=r(t)\,e^{\lambda_jt}\ (\deg r=l)\ \text{ if }\lambda_j\ne\lambda,\qquad g=\frac{t^{l+1}e^{\lambda t}}{l+1}\ \text{ if }\lambda_j=\lambda .
\]
With $v_p$ their sum, $\partial_t\big(e^{-\lambda t}(v-v_p)\big)=0$.

A distribution $V$ on $J'\times I'$ with $\partial_tV=0$ is $c\otimes1$. Indeed, fix $\chi\in C_c^\infty(I')$ with $\int\chi=1$. For a test function $\varphi$ put $\Phi(x)=\int\varphi(x,t)\dd t$ and
\[
\psi(x,t)=\int_{-\infty}^t\big(\varphi(x,s)-\chi(s)\Phi(x)\big)\dd s .
\]
Since $I'$ is an interval, $\psi$ has compact support, and $\varphi=\chi\otimes\Phi+\partial_t\psi$. Hence $V(\varphi)=V(\chi\otimes\Phi)$, which depends only on $\Phi$. So $v=v_p+c\otimes e^{\lambda t}$, completing the induction.

\step{7}{The ODE in $x$ makes the coefficients exponential-polynomials}
Applying $q(\partial_x)$:
\[
\sum_{j,l}q(\partial_x)c_{jl}\otimes t^le^{\lambda_jt}=0 .
\]
The functions $g_a=t^le^{\lambda_jt}$ are independent on $I'$, so there are test functions $\psi_b$ with $\int g_a\psi_b=\delta_{ab}$. Testing against $\varphi(x)\psi_b(t)$ gives $q(\partial_x)c_{jl}=0$. The one-variable case of Step~6 gives $c_{jl}\in\EP(J')$. So $w$ restricted to $Q_0$ is the restriction of an exponential-polynomial $E_{Q_0}$ on $\R^2$.

\step{8}{Connectedness glues the rectangles}
Fix $Q_0$ and $E=E_{Q_0}$. Let $A=\{y\in\Omega:\ w=E\ \text{near }y\}$. It is open and nonempty. It is closed in $\Omega$: if $y\in\bar A$, take a rectangle $Q_y\ni y$. On the open set $Q_y\cap A$ we have $E_{Q_y}=w=E$, so by Step~2 the exponential-polynomial $E_{Q_y}-E$ vanishes identically, and $w=E$ near $y$. By connectedness $A=\Omega$, so $w=E|_\Omega\in\EP(\Omega)$.

\whyworks{finite dimension turns differentiation into two commuting matrices; their eigenvalues are the exponents, and rectangles plus connectedness globalise the local form.}
\end{proof}
\uses{A1.}

\begin{caution}[title={Why rectangles in Step~6}]
On a U-shaped region, $v=\phi(x)\mathbf 1_{t>2}$ (with $\phi=0$ for $x<1$, $\phi=1$ for $x>2$) is smooth and satisfies $\partial_tv=0$, but it is not of the form $c(x)\otimes1$. So the local form is proved on rectangles and glued in Step~8. This $v$ is not a counterexample to A2: its derivative space is infinite-dimensional.
\end{caution}

%=====================================================================
\section{The identification theorems}

\begin{dimcard}{A3}{C3-T1}{Identification criterion}{\proved}{thm:ident}
Let $\Omega$ be connected, $u\in\mathcal D'(\Omega)$, $P(\partial)u=0$ with $P=c\prod_iP_i^{e_i}$ nonconstant, and
\[
w_i=\Big(\frac{P}{P_i}\Big)(\partial)u,\qquad w_i'=\Big(\frac{P}{P_i^{e_i}}\Big)(\partial)u .
\]
\begin{enumerate}[label=(\alph*)]
\item $\Ann(u)=(P)\iff w_i\notin\EP(\Omega)$ for every $i$.
\item $\Var(\Ann u)=\Var(P)\iff w_i'\notin\EP(\Omega)$ for every $i$.
\item If $P$ is squarefree, the two conditions coincide.
\item If $P$ is monic in $\tau$ of degree $m$: $\Ann(u)=(P)$ iff the only $(b_0,\dots,b_{m-1})\in\C[\xi]^m$ with $\sum_jb_j(\partial_x)\partial_t^ju=0$ is zero.
\item In general $\Ann(u)=G\cdot J$ with $G\mid P$, $\Var(J)$ finite, and $G\sim P\iff\Ann(u)=(P)$.
\end{enumerate}
\end{dimcard}

\begin{proofidea}
Write $\Ann(u)=G\cdot J$. If $G$ misses a copy of some $P_i$, then $w_i$ is killed by a power of the ideal of the finite set $\Var(J)$, so it is an exponential-polynomial. Conversely, an exponential-polynomial $w_i$ is killed by some $h$ coprime to $P_i$, and $h\cdot P/P_i$ is an annihilator of $u$ that $P$ does not divide.
\end{proofidea}

\begin{proof}
Let $I=\Ann(u)$. It contains $P$. By A1, $I=G\cdot J$ with $G=\Gcd I$ dividing $P$, so
\[
G=c'\prod_iP_i^{f_i},\qquad 0\le f_i\le e_i .
\]

\step{1}{(e) $G\sim P$ exactly when $I=(P)$}
If $G\sim P$, then $I=PJ\subseteq(P)\subseteq I$. If $I=(P)$, then $P$ divides every element of $I$ and is itself in $I$, so it is the gcd.

\step{2}{(a), first direction: if $I\ne(P)$, some $w_i$ is an exponential-polynomial}
If $I\ne(P)$, then $G\not\sim P$, so some $f_i<e_i$, and $P/P_i=G\,H$ for some $H\in\ring$. By A1(iii), $J\supseteq\mathfrak n^N$ with $\mathfrak n$ the ideal of $\Var(J)$. For $h\in\mathfrak n^N$:
\[
h\cdot\frac{P}{P_i}=G\,(Hh)\in G\,J=I,
\qquad\text{so}\qquad h(\partial)\,w_i=0 .
\]
Hence $\Ann(w_i)\supseteq\mathfrak n^N$, its zero set lies in the finite $\Var(J)$, and $w_i\in\EP$ by A2.

\step{3}{(a), second direction: an exponential-polynomial $w_i$ gives a new annihilator}
Suppose $w_i\in\EP$. Then $\Var(\Ann w_i)$ is finite, so by A1(F) some $h\in\Ann(w_i)$ has $P_i\nmid h$. Put
\[
Q=h\cdot\frac{P}{P_i}.
\]
Then $Q(\partial)u=h(\partial)w_i=0$, so $Q\in I$. If $Q=PA$, cancelling $P/P_i$ in the domain $\ring$ gives $h=P_iA$, a contradiction. So $Q\in I\setminus(P)$.

\step{4}{(b) The variety sees only whether each $f_i\ge1$}
We have $\Var(I)=\Var(G)\cup\Var(J)$, with $\Var(J)$ finite. If $f_i\ge1$, then $\Var(P_i)\subseteq\Var(G)\subseteq\Var(I)$. If $f_i=0$, then $\Var(P_i)\cap\Var(I)$ lies in the finite set
\[
\bigcup_{f_j\ge1}\big(\Var(P_i)\cap\Var(P_j)\big)\ \cup\ \Var(J),
\]
while $\Var(P_i)$ is infinite, so $\Var(P_i)\not\subseteq\Var(I)$. Hence $\Var(I)=\Var(P)$ iff all $f_i\ge1$. Finally, $f_i=0$ iff $w_i'\in\EP$: one direction is Step~2 with $P/P_i^{e_i}$ in place of $P/P_i$; for the other, an $h\in\Ann(w_i')$ with $P_i\nmid h$ gives $G\mid h\,P/P_i^{e_i}$, and $P_i$ divides neither factor, so $f_i=0$.

\step{5}{(c) and (d)}
For $e_i=1$, $w_i=w_i'$, which is (c). For (d), divide by the monic $P$ in $\C[\xi][\tau]$:
\[
Q=AP+\sum_{j<m}b_j(\xi)\,\tau^j\quad\text{uniquely}.
\]
Since $A(\partial)P(\partial)u=0$, $Q\in I$ iff $\sum_jb_j(\partial_x)\partial_t^ju=0$, and $Q\in(P)$ iff all $b_j=0$.

\whyworks{the gcd of the annihilator records the infinite part of $u$; any factor it misses leaves behind a function killed by a zero-dimensional ideal, that is, an exponential-polynomial.}
\end{proof}
\uses{A1, A2.}

\begin{example}{The proof on a wave}
Let $P=\tau^2-\xi^2$ and $u=e^{x-t}+e^{-(x+t)^2}$. For the factor $P_i=\tau+\xi$:
\[
w_i=(\partial_t-\partial_x)u=-2e^{x-t}\in\EP .
\]
Step~3 takes $h=\xi-1$, which kills $e^{x-t}$ and is not a multiple of $\tau+\xi$, and produces
\[
Q=(\xi-1)(\tau-\xi)\in\Ann(u)\setminus(P).
\]
\end{example}

%---------------------------------------------------------------------
\begin{dimcard}{A4}{C3-T6}{The intrinsic operator}{\proved}{thm:intrinsic}
Let $0\ne P\in\Ann(u)$, $P=c\prod_iP_i^{e_i}$, and
\[
S_i=\Big\{k\in\{0,\dots,e_i\}:\ \Big(\frac{P}{P_i^k}\Big)(\partial)u\in\EP\Big\}.
\]
Then $G_u=\Gcd\Ann(u)\sim\prod_iP_i^{f_i}$ with $f_i=e_i-\max S_i$. Moreover $\Ann(u)$ is principal iff $G_u(\partial)u=0$, and then $\Ann(u)=(G_u)$.
\end{dimcard}

\begin{proofidea}
The multiplicity $f_i$ of $P_i$ in $G_u$ is detected by stripping copies of $P_i$ from $P$ one at a time: as long as fewer than $e_i-f_i$ copies are removed the leftover is an exponential-polynomial, and as soon as more are removed it is not.
\end{proofidea}

\begin{proof}
\step{1}{The key equivalence}
By A1, $\Ann(u)=G_uJ_u$ with $G_u\mid P$. We show, for $1\le k\le e_i$,
\[
f_i\le e_i-k\iff\Big(\frac{P}{P_i^k}\Big)(\partial)u\in\EP .
\]
($\Rightarrow$) Then $G_u\mid P/P_i^k$; write $P/P_i^k=G_uH$. As in A3, Step~2, $h\cdot P/P_i^k\in\Ann(u)$ for every $h$ in a power of the ideal of $\Var(J_u)$, so the function is an exponential-polynomial.

($\Leftarrow$) If it is an exponential-polynomial, pick $h$ in its annihilator with $P_i\nmid h$ (A1(F)). Then $G_u$ divides $h\cdot P/P_i^k$. Comparing multiplicities of $P_i$ gives $f_i\le e_i-k$.

\step{2}{Reading off $f_i$}
The equivalence holds trivially for $k=0$. So $S_i=\{0,1,\dots,e_i-f_i\}$, and $\max S_i=e_i-f_i$. (The set is an interval because $\EP$ is stable under derivatives:
\[
\Big(\frac{P}{P_i^{k-1}}\Big)(\partial)u=P_i(\partial)\Big[\Big(\frac{P}{P_i^k}\Big)(\partial)u\Big].)
\]

\step{3}{Principality}
If $G_u(\partial)u=0$, then $(G_u)\subseteq\Ann(u)=G_uJ_u\subseteq(G_u)$. Conversely, if $\Ann(u)=(Q)$, then $Q$ is the gcd: $Q\sim G_u$, and $G_u(\partial)u=0$.

\whyworks{each copy of $P_i$ that $u$ does not need leaves an exponential-polynomial behind; counting them gives the exponent.}
\end{proof}
\uses{A1, A2, A3.}

%---------------------------------------------------------------------
\begin{dimcard}{A5}{P4-T2}{Spectral theorem for the cyclic module}{\proved}{thm:spec}
Let $W=\ring u\cong\ring/I$ with $I=\Ann(u)=G\cdot J$ as in A1. A \emph{joint eigenvector} for $z\in\C^2$ is $w\in W$, $w\ne0$, with $\mathfrak m_z(\partial)w=0$; $\sigma_p(u)$ is the set of $z$ having one.
\begin{enumerate}[label=(\roman*)]
\item Every joint eigenvector for $z$ is $c\,e^{\langle z,\cdot\rangle}$ with $c\ne0$, and $\sigma_p(u)=\Var(J)$.
\item $W\cap\EP(\Omega)=\ring\,G(\partial)u\cong\ring/J$, of finite dimension $\dimc\ring/J$.
\item $\Ann(u)$ is principal $\iff J=\ring\iff\sigma_p(u)=\emptyset\iff W\cap\EP=0$.
\item Hence, if $P(\partial)u=0$: $\Ann(u)=(P)\iff\sigma_p(u)=\emptyset$ and $G\sim P$.
\end{enumerate}
\end{dimcard}

\begin{proofidea}
Translate statements about the quotient ring $\ring/I$ (maximal ideals, socle) into statements about $W$. The only analytic input is the EP test A2.
\end{proofidea}

\begin{proof}
Two small facts are used.
\begin{itemize}
\item \textbf{(LG)} If $G\mid(\tau-\tau_z)Q$ and $G\mid(\xi-\xi_z)Q$, then $G\mid Q$: in a unique factorisation domain the gcd of the two products is $Q$.
\item \textbf{(LS)} If $M\ne0$ is a module with $\mathfrak m_z^NM=0$, some $v\ne0$ in $M$ has $\mathfrak m_zv=0$: take $y\ne0$ and $k$ maximal with $\mathfrak m_z^ky\ne0$.
\end{itemize}

\step{1}{Joint eigenvectors are single exponentials}
If $w=Q(\partial)u\ne0$ and $\mathfrak m_z^N(\partial)w=0$, then $\Var(\Ann w)\subseteq\{z\}$, so $w\in\EP$ and by A2 $w=q\,e^{\langle z,\cdot\rangle}$. The eigenvector condition gives $\partial_tq=\partial_xq=0$: $q$ is a constant.

\step{2}{$\sigma_p(u)\subseteq\Var(J)$}
Let $Q(\partial)u$ be a joint eigenvector at $z$. Then $(\tau-\tau_z)Q$ and $(\xi-\xi_z)Q$ lie in $I\subseteq(G)$, so $G\mid Q$ by (LG). Write $Q=GQ'$. Cancelling $G$:
\[
(\tau-\tau_z)Q'\in J,\qquad (\xi-\xi_z)Q'\in J,\qquad Q'\notin J .
\]
So $\mathfrak m_z\subseteq(J:Q')\ne\ring$. As $\mathfrak m_z$ is maximal, $(J:Q')=\mathfrak m_z\supseteq J$, and $z\in\Var(J)$.

\step{3}{$\Var(J)\subseteq\sigma_p(u)$}
By A1(iii), $J\supseteq\prod_{z'\in\Var(J)}\mathfrak m_{z'}^N$, and by the Chinese remainder theorem
\[
\ring/J\cong\prod_{z'\in\Var(J)}\ring/\big(J+\mathfrak m_{z'}^N\big).
\]
The factor at $z$ is nonzero and killed by $\mathfrak m_z^N$. By (LS) it contains $v\ne0$ with $\mathfrak m_zv=0$. Through the isomorphism of Step~4 this gives a joint eigenvector $(vG)(\partial)u\ne0$.

\step{4}{The exponential-polynomial part is $\ring G(\partial)u$}
Since $\ring$ is a domain,
\[
\Ann\big(G(\partial)u\big)=(I:G)=\{Q:\ QG\in GJ\}=J .
\]
$\Var(J)$ is finite, so $G(\partial)u\in\EP$ and $\ring G(\partial)u\subseteq\EP$. Conversely, let $w=Q(\partial)u\in\EP$. For each prime factor $P_i$ of $G$, with multiplicity $f_i$, A1(F) gives $h\in\Ann(w)$ with $P_i\nmid h$. Then $hQ\in I\subseteq(G)$, so $P_i^{f_i}\mid hQ$ and $P_i^{f_i}\mid Q$. Hence $G\mid Q$ and $w\in\ring G(\partial)u$. The map $Q\mapsto(QG)(\partial)u$ is onto with kernel $J$, so the space is $\cong\ring/J$.

\step{5}{(iii) and (iv)}
$J=\ring\iff\Var(J)=\emptyset$ (A1(iii)) $\iff\sigma_p(u)=\emptyset$ (Steps~2--3). Also $J=\ring\iff I=(G)$ is principal, and $W\cap\EP\cong\ring/J=0\iff J=\ring$. Finally $\Ann(u)=(P)$ iff $I$ is principal and $G\sim P$.

\whyworks{the finite part $J$ of the annihilator is exactly the module of exponentials inside $\ring u$; its points are the eigenvalues.}
\end{proof}
\uses{A1, A2.}

\begin{caution}[title={An empty spectrum alone is not enough}]
Take $u=e^{x^2}$ and $P=\tau^2$. For $Q=Q_0(\xi)+\tau Q_1$ one has $Q(\partial)u=Q_0(\partial_x)e^{x^2}$, which vanishes only if $Q_0=0$. So $\Ann(u)=(\tau)$ is principal and $\sigma_p(u)=\emptyset$, but $\Ann(u)\ne(\tau^2)$. The spectrum sees $J_u$, not $G_u$; condition (iv) needs both.
\end{caution}

%=====================================================================
\section{Finite data}

\begin{dimcard}{A6}{C3-T2}{B\'ezout threshold and jet threshold}{\proved}{thm:bezout}
\textbf{(C3-T2)} Let $u=\sum_{k=1}^Mc_k\,e^{\langle z_k,\cdot\rangle}$ with distinct $z_k\in\Var(P)$ and $c_k\ne0$, $P$ squarefree, $n\ge\deg P$. If
\[
\big|\{z_k\}\cap\Var(P_i)\big|>n\deg P_i\qquad\text{for every }i,
\]
then $\Ann(u)\cap\ring_{\le n}=P\cdot\ring_{\le n-\deg P}$. The threshold cannot be lowered uniformly.

\phantomsection\label{dim:C3-T3}\textbf{(C3-T3)} Let $0\ne u\in\EP(\Omega)$, $d=\dimc\ring u$, $p_0\in\Omega$.
\begin{enumerate}[label=(\roman*)]
\item $\Ann(u)\cap\ring_{\le n}=\{Q\in\ring_{\le n}:\ \partial^\beta(Q(\partial)u)(p_0)=0\ \forall|\beta|\le d-1\}$.
\item $\Ann(u)$ is generated in degree $\le d$, so the jet of order $2d-1$ determines it.
\item The order $2d-1$ is optimal in the worst case.
\end{enumerate}
\end{dimcard}

\begin{proofidea}
For modes: an annihilator vanishes at every exponent, and B\'ezout bounds how many exponents a curve that does not contain $\Var(P_i)$ can pass through. For jets: inside the $d$-dimensional space $\ring u$, the subspaces of elements vanishing to order $k$ at $p_0$ shrink strictly, so after $d-1$ steps they are zero.
\end{proofidea}

\begin{proof}
\step{1}{Annihilators of a sum of modes}
By A2(a)(b), $Q(\partial)u=\sum_kc_kQ(z_k)\,e^{\langle z_k,\cdot\rangle}$, which is zero iff $Q(z_k)=0$ for all $k$.

\step{2}{B\'ezout forces every $P_i$ to divide $Q$}
Let $Q\in\Ann(u)$ with $\deg Q\le n$. If $Q=0$ there is nothing to prove; a nonzero constant cannot kill $u\ne0$. Suppose $P_i\nmid Q$. Then $Q$ and $P_i$ are coprime, and B\'ezout's theorem (the homogenisations have no common component, since the line at infinity divides neither) gives
\[
\big|\Var(Q)\cap\Var(P_i)\big|\le\deg Q\cdot\deg P_i\le n\deg P_i .
\]
But by Step~1, $\{z_k\}\cap\Var(P_i)\subseteq\Var(Q)\cap\Var(P_i)$ has more points. So $P_i\mid Q$ for every $i$, and since the $P_i$ are distinct primes, $P\mid Q$.

\step{3}{The threshold is sharp}
Take $Z_n=\{(j^2,\pm j):\ j=1,\dots,n\}$, $2n$ points on $\tau=\xi^2$, and
\[
Q_n=\prod_{j=1}^n(\tau-j^2).
\]
It has degree $n$, vanishes on $Z_n$, and is not a multiple of $\tau-\xi^2$ because it does not involve $\xi$.

\step{4}{(i) Vanishing subspaces shrink strictly}
Let $W=\ring u$ and
\[
K_k=\{w\in W:\ \partial^\beta w(p_0)=0\ \text{for all }|\beta|\le k\}.
\]
If $K_k=K_{k+1}$, then $K_k$ is stable under derivatives (if $w\in K_{k+1}$ then $\partial w\in K_k$), so every $w\in K_k$ has all derivatives zero at $p_0$; being analytic on a connected set, $w=0$. So the dimensions drop strictly until they reach $0$. Since $\dim K_0=d-1$, we get $K_{d-1}=0$. For $Q\in\ring_{\le n}$, $Q(\partial)u\in W$ vanishes iff it lies in $K_{d-1}$. Any upper bound $d'\ge d$ works as well.

\step{5}{(ii) Generators of low degree}
$\ring/\Ann(u)\cong W$ has dimension $d$. In a graded monomial order, the standard monomials form a set of $d$ elements closed under division. A standard monomial $\tau^a\xi^b$ has $(a+1)(b+1)\ge a+b+1$ divisors, all standard, so $a+b+1\le d$. The minimal generators of the initial ideal therefore have degree at most $d$, and so do the elements of the reduced Gr\"obner basis, which generate $\Ann(u)$. Apply (i) with $n=d$: the jet of order $2d-1$ suffices.

\step{6}{(iii) Optimality}
Let $u=\sum_{j=1}^de^{jx}$, so $d=\dimc\ring u$, and $n=d$. Then $\Ann(u)\cap\C[\xi]_{\le d}=\C\cdot Q_0$ with $Q_0=\prod_j(\xi-j)$. Put $m_k=\partial_x^ku(p_0)$. With only $|\beta|\le d-2$ the test of (i) is the Hankel system
\[
\sum_{i=0}^dq_i\,m_{k+i}=0,\qquad k=0,\dots,d-2,
\]
which has $d-1$ equations in $d+1$ unknowns, so a kernel $N$ of dimension at least $2$. Take $Q_1\in N\setminus\C Q_0$ and $Q'=Q_0+sQ_1$ with $s$ small: it has degree $d$ and $d$ simple roots $\mu_j$. Choose $w_j$ with $\sum_jw_j\mu_j^k=m_k$ for $k<d$ (Vandermonde) and set
\[
u'=\sum_jw_j\,e^{\mu_j(x-x_0)}.
\]
The moment sequences of $u$ and $u'$ satisfy the same recursion given by $Q'$ and agree for $k\le d-1$, hence for $k\le2d-2$. So $u$ and $u'$ have the same jet of order $2d-2$, but $Q'\in\Ann(u')\setminus\Ann(u)$.

\whyworks{finitely many exponentials are finitely many points; B\'ezout counts how many points a curve must miss, and the strict drop of vanishing subspaces counts how many derivatives are needed.}
\end{proof}
\uses{A1, A2.}

%---------------------------------------------------------------------
\begin{dimcard}{A7}{C3-P1}{No finite certificate}{\proved{} (jets) / \cf{} (functionals)}{thm:nocert}
Let $\Omega$ be connected, $p_0\in\Omega$, $P$ nonconstant, $u$ smooth with $P(\partial)u=0$, $K\ge0$, and $E_P=\ker P\cap\EP(\Omega)$.
\begin{enumerate}[label=(\roman*)]
\item $\bigcap_{v\in E_P}\Ann(v)=(P)$.
\item There is $v\in E_P$ with the same jet of order $K$ at $p_0$ as $u$; in particular $\Ann(v)\supsetneq(P)$.
\end{enumerate}
On convex $\Omega$ the same holds for any finitely many continuous linear functionals (Malgrange); on a non-convex region it can fail.
\end{dimcard}

\begin{proofidea}
By duality, the possible jets of solutions are cut out by the multiples of $P$. Exponential-polynomial solutions already cut out the same space, because for every $Q$ that is not a multiple of $P$ one builds a solution $s^{e-1}e^{\langle z,\cdot\rangle}$ that $Q$ does not kill.
\end{proofidea}

\begin{proof}
Pair an operator $Q$ of degree at most $K$ with a jet by
\[
\langle Q,J_Kw\rangle=\big(Q(\partial)w\big)(p_0).
\]

\step{1}{Vanishing on all of $E_P$ means annihilating all of $E_P$}
If $(Q(\partial)v)(p_0)=0$ for every $v\in E_P$, then, since $E_P$ is stable under derivatives, all derivatives of the exponential-polynomial $Q(\partial)v$ vanish at $p_0$. It is entire with zero Taylor series, so it is $0$. Hence $Q\in\bigcap_{v\in E_P}\Ann(v)$.

\step{2}{(i) A non-multiple of $P$ misses some exponential solution}
Let $Q\ne0$ with $P\nmid Q$. Some $i$ has $f=v_{P_i}(Q)<e_i$; write $Q=P_i^fQ'$ with $P_i\nmid Q'$. The sets $\Var(P_i)\cap\Var(\nabla P_i)$ and $\Var(P_i)\cap\Var(Q')$ are finite (A1(V)), while $\Var(P_i)$ is infinite. Pick $z\in\Var(P_i)$ with
\[
\nabla P_i(z)\ne0,\qquad Q'(z)\ne0,
\]
and a coordinate $s\in\{t,x\}$ with $\partial_sP_i(z)\ne0$. On polynomials in $s$, $P_i(z+\partial)$ lowers the degree by exactly one and $Q'(z+\partial)$ preserves it. Let $v=s^{e_i-1}e^{\langle z,\cdot\rangle}$. Then $P(z+\partial)s^{e_i-1}=0$, so $v\in E_P$, while $Q(z+\partial)s^{e_i-1}$ has degree $e_i-1-f\ge0$ and is nonzero. So $Q\notin\Ann(v)$.

\step{3}{(ii) Every jet of a solution is a jet of an exponential solution}
By Steps~1--2, every $Q$ of degree $\le K$ orthogonal to the jets of $E_P$ is a multiple of $P$, hence orthogonal to the jet of $u$. In finite dimension this means
\[
J_K(u)\in J_K(E_P),
\]
so a suitable $v\in E_P$ exists. Finally $\Ann(v)\supsetneq(P)$: if $v=0$ it is $\ring$; otherwise $\Var(\Ann v)$ is finite while $\Var(P)$ is infinite.

\step{4}{General functionals and convexity}
For general continuous functionals the analogue on convex $\Omega$ is Malgrange's approximation theorem (Ann.\ Inst.\ Fourier 6, Chap.~I \S2 Thm~2) \cf. Without convexity it fails: on a U-shaped region take $P=\tau$, $u=\phi(x)\mathbf 1_{t>2}$, and
\[
\ell(w)=w(5/2,5/2)-w(1/2,5/2).
\]
Then $\ell(u)=1$ while $\ell(v)=0$ for every $v\in\ker\partial_t$ with $\Ann(v)\ne(\tau)$.

\whyworks{finite data see a finite-dimensional shadow of $\ker P$, and exponential solutions already fill every such shadow.}
\end{proof}
\uses{A1, A2.}

%---------------------------------------------------------------------
\begin{dimcard}{A8}{P4-T4}{Decidability with a bound}{\proved}{thm:decide}
In the model M2$^+$ (exact jets at $p_0$ in a computable field $K$, plus a zero-test oracle):
\begin{enumerate}[label=(\roman*)]
\item for every finite-dimensional $\mathcal L\subseteq\ring$ with a basis over $K$, $\Ann(u)\cap\mathcal L$ is computable;
\item ``$\Ann(u)\ne(P)$'' is semi-decidable;
\item if an integer $B\ge\dimc(\ring u\cap\EP)$ is given, then with $N=B+\deg P-1$,
\[
\Ann(u)=(P)\iff\Ann(u)\cap\ring_{\le N}=(P)\cap\ring_{\le N},
\]
which is decidable. The test ``$Q(\partial)u\in\EP$?'' is then decidable, and $G_u$ is computable.
\end{enumerate}
\end{dimcard}

\begin{proofidea}
Linear algebra on jets gives a decreasing chain of candidate spaces that stabilises at the true annihilators; the zero oracle confirms when it has stabilised. If identification fails, a witness of small degree exists, and the bound $B$ says how small.
\end{proofidea}

\begin{proof}
\step{1}{(i) A decreasing chain of linear spaces}
Let
\[
A_k=\{Q\in\mathcal L:\ \partial^\beta\big(Q(\partial)u\big)(p_0)=0\ \text{for all }|\beta|\le k\}.
\]
It is the kernel of a matrix with entries $\partial^{\beta+\alpha}u(p_0)\in K$, so it has a computable basis. The chain $A_0\supseteq A_1\supseteq\dots\supseteq\Ann(u)\cap\mathcal L$ stabilises, and its limit is $\Ann(u)\cap\mathcal L$, because $Q(\partial)u$ is analytic with zero Taylor series at $p_0$. Algorithm: for $k=0,1,\dots$ compute a basis of $A_k$ and ask the oracle about each element; stop when all answers are ``zero''. It stops only when $A_k\subseteq\Ann(u)$, so it is correct, and it stops at the stabilisation index.

\step{2}{(ii) Semi-decision}
For $n=\deg P,\deg P+1,\dots$ compute $\Ann(u)\cap\ring_{\le n}$ by (i) and compare its dimension with that of $P\cdot\ring_{\le n-\deg P}$. If $\Ann(u)\ne(P)$, a witness $Q$ exists and the loop stops at $n=\max(\deg Q,\deg P)$.

\step{3}{(iii) A witness of degree at most $N$}
Suppose $\Ann(u)\ne(P)$. By A3, some $w_i=(P/P_i)(\partial)u$ is an exponential-polynomial in $\ring u$, so
\[
d_i=\dimc\ring w_i\le B .
\]
By A6(ii), $\Ann(w_i)$ is generated in degree $\le d_i$, and not all generators are divisible by $P_i$ (A1(F)). Pick such a generator $h$, so $\deg h\le B$. Then
\[
Q=h\cdot\frac{P}{P_i}\in\Ann(u)\setminus(P),\qquad \deg Q\le B+\deg P-1=N .
\]
The comparison in degree $N$ is computable by (i).

\step{4}{The EP test from the bound}
Let $w=Q(\partial)u$; its oracles follow from those of $u$. If $w\in\EP$, then $\dimc\ring w\le B$, so the minimal polynomials of $\partial_t$ and $\partial_x$ on $\ring w$ have degree $\le B$: there are nonzero $a(\tau),b(\xi)$ of degree $\le B$ in $\Ann(w)$. Conversely, such $a,b$ give a finite zero set and $w\in\EP$ (A2). Both searches are computable by (i). With this test, A4 computes $G_u$.

\whyworks{a failure of identification is always witnessed in bounded degree, and the bound $B$ makes that degree explicit.}
\end{proof}
\uses{A1, A2, A3, A4, A6.}

%---------------------------------------------------------------------
\begin{dimcard}{A9}{P4-T5}{Black-box impossibility}{\proved}{thm:blackbox}
Let $P\in K[\tau,\xi]$ be nonconstant ($K$ the algebraic numbers), $\Omega$ connected, $p_0\in\Omega$ with coordinates in $K\cap\R$, and $u$ real-analytic with $P(\partial)u=0$, jets in $K$, and $\Ann(u)=(P)$ (such $u$ exists for every $P$, \led{P4-L7}). For every order $k_0$ and every finite set $\mathcal Q\subset K[\tau,\xi]$ there is an exponential-polynomial $v\in\ker P$ with
\begin{itemize}
\item the same jet of order $k_0$ at $p_0$ as $u$;
\item for every $Q\in\mathcal Q$: $Q(\partial)v=0\iff Q(\partial)u=0$;
\item $\Ann(v)\supsetneq(P)$.
\end{itemize}
Hence no algorithm in the model M2$^+$ decides ``$\Ann(\cdot)=(P)$'' correctly on all inputs.
\end{dimcard}

\begin{proofidea}
Run a correct algorithm on $u$: it halts after finitely many questions. An adversary builds $v$ among the exponential-polynomial solutions so that it gives the same answer to every one of those questions. First match the jet, then, for each $Q$ with $Q(\partial)u\ne0$, add a solution with zero jet that $Q$ does not kill, with generic coefficients.
\end{proofidea}

\begin{proof}
\step{1}{Witness points with algebraic coordinates}
Each $\Var(P_i)$ has infinitely many points with coordinates in $K$: roots of $P_i(\cdot,\xi_0)$ for $\xi_0\in K$, or $K\times\{\xi_0\}$ if $P_i$ does not involve $\tau$. For such $z$ with $\nabla P_i(z)\ne0$, pick a coordinate $s_z$ with $\partial_{s}P_i(z)\ne0$ and put
\[
g_{i,z,j}=s_z^{\,j}\,e^{\langle z,\cdot-p_0\rangle},\qquad j<e_i .
\]
As in A7, Step~2, $P(z+\partial)s_z^j=0$, so $g_{i,z,j}\in\ker P$.

\step{2}{Matching the jet}
Let $V_{k_0}$ be the span of the jets of order $k_0$ of all $g_{i,z,j}$. With the pairing of A7,
\[
V_{k_0}^\perp=(P)\cap\ring_{\le k_0};
\]
the inclusion ``$\subseteq$'' is A7, Step~2, with points chosen in $K^2$. Since $\langle PA,J_{k_0}u\rangle=0$, $J_{k_0}(u)\in V_{k_0}$. A linear system with coefficients in $K$ that is solvable over $\C$ is solvable over $K$, so some combination $v_1$ of the $g$'s has the jet of $u$.

\step{3}{A zero-jet solution for each test}
For $Q\in\mathcal Q\setminus(P)$ (so $Q(\partial)u\ne0$), pick $i$ with $f=v_{P_i}(Q)<e_i$, $Q=P_i^fQ'$, and more than $\dim\ring_{\le k_0}$ points $z\in\Var(P_i)\cap K^2$ with $\nabla P_i(z)\ne0$ and $Q'(z)\ne0$. The map
\[
c\ \mapsto\ J_{k_0}\Big(\sum_zc_z\,s_z^{e_i-1}e^{\langle z,\cdot-p_0\rangle}\Big)
\]
goes from a larger space to a smaller one, so it has a nonzero kernel vector $c$. Let $e_Q$ be the corresponding sum: its jet is $0$, while $Q(\partial)e_Q\ne0$ by A2(b).

\step{4}{A generic combination answers like $u$}
Let $v_\lambda=v_1+\sum_Q\lambda_Q\,e_Q$. It has the jet of $u$ for every $\lambda$. For each $Q$ with $Q(\partial)u\ne0$, the set $\{\lambda:\ Q(\partial)v_\lambda=0\}$ is a proper affine subspace. Finitely many proper subspaces do not cover the space, so some $\lambda$ avoids all of them. For $Q\in(P)$ both answers are ``zero''. So $v=v_\lambda$ answers every question as $u$ does, and $v\in\EP\cap\ker P$ gives $\Ann(v)\supsetneq(P)$.

\step{5}{Conclusion}
The run of the algorithm on $u$ depends only on the answers it receives. They coincide for $v$, so it outputs ``$\Ann=(P)$'' on $v$ as well, which is wrong.

\whyworks{a finite conversation with a black box can always be imitated by a finite sum of exponentials, and such sums never identify.}
\end{proof}
\uses{A2, A7.}
